\chapter{Relational Persistence, Vector Indexing and Hybrid Search}
\label{chap:storage_search}

\section{SQLite Architecture and Database Pragmas}
Conversational state, reasoning traces, and metadata are persisted locally in \texttt{conversations.db}. Finch.ai configures SQLite with high-concurrency pragmas upon connection initialization:

\begin{lstlisting}[language=SQL, caption={SQLite Production Connection Pragmas}]
PRAGMA journal_mode = WAL;
PRAGMA synchronous = NORMAL;
PRAGMA foreign_keys = ON;
PRAGMA busy_timeout = 5000;
\end{lstlisting}

Write-Ahead Logging (WAL) permits concurrent reader processes without locking the database file, allowing real-time CLI queries while background summarizers and active chats commit new turns.

\section{Relational Schema}
The persistence model is composed of two primary relational tables and two specialized virtual tables.

\subsection{The \texttt{sessions} Table}
Tracks conversational sessions, lifecycle timestamps, summary texts, and token totals:
\begin{lstlisting}[language=SQL]
CREATE TABLE IF NOT EXISTS sessions (
    id TEXT PRIMARY KEY,
    title TEXT,
    created_at TIMESTAMP DEFAULT CURRENT_TIMESTAMP,
    updated_at TIMESTAMP DEFAULT CURRENT_TIMESTAMP,
    summary TEXT,
    total_tokens INTEGER DEFAULT 0,
    model TEXT,
    provider TEXT
);
\end{lstlisting}

\subsection{The \texttt{messages} Table}
Maintains the complete sequential dialogue history with foreign-key cascades:
\begin{lstlisting}[language=SQL]
CREATE TABLE IF NOT EXISTS messages (
    id TEXT PRIMARY KEY,
    session_id TEXT NOT NULL REFERENCES sessions(id) ON DELETE CASCADE,
    role TEXT NOT NULL,
    content TEXT NOT NULL,
    created_at TIMESTAMP DEFAULT CURRENT_TIMESTAMP,
    model TEXT,
    prompt_tokens INTEGER DEFAULT 0,
    eval_tokens INTEGER DEFAULT 0,
    eval_duration_ms REAL DEFAULT 0.0,
    tokens_per_second REAL DEFAULT 0.0,
    ttft_ms REAL DEFAULT 0.0,
    reasoning_content TEXT,
    tool_calls TEXT,
    tool_call_id TEXT,
    compressed INTEGER DEFAULT 0
);
\end{lstlisting}

\section{Sub-Millisecond Lexical Search with FTS5}
To deliver instantaneous recall of technical terms, code snippets, dates, and filenames, Finch.ai provisions an SQLite FTS5 full-text search virtual table:

\begin{lstlisting}[language=SQL]
CREATE VIRTUAL TABLE IF NOT EXISTS messages_fts USING fts5(
    message_id UNINDEXED,
    session_id UNINDEXED,
    role,
    content,
    tokenize = 'porter unicode61'
);
\end{lstlisting}

\subsection{Trigger-Driven Synchronization}
To eliminate indexing lag and maintain ACID consistency, three database triggers synchronize \texttt{messages\_fts} synchronously on write operations:

\begin{lstlisting}[language=SQL]
CREATE TRIGGER IF NOT EXISTS trg_messages_ai AFTER INSERT ON messages BEGIN
    INSERT INTO messages_fts(message_id, session_id, role, content)
    VALUES (new.id, new.session_id, new.role, new.content);
END;

CREATE TRIGGER IF NOT EXISTS trg_messages_ad AFTER DELETE ON messages BEGIN
    DELETE FROM messages_fts WHERE message_id = old.id;
END;

CREATE TRIGGER IF NOT EXISTS trg_messages_au AFTER UPDATE ON messages BEGIN
    DELETE FROM messages_fts WHERE message_id = old.id;
    INSERT INTO messages_fts(message_id, session_id, role, content)
    VALUES (new.id, new.session_id, new.role, new.content);
END;
\end{lstlisting}

\section{Dense Semantic Vector Search with \texttt{sqlite-vec}}
While FTS5 excels at exact keyword matching (e.g., \texttt{def calculate\_hash()}), natural language queries frequently use conceptual, non-overlapping vocabulary (e.g., "how did we handle caching in python?").

Finch.ai uses the local embedding model \texttt{nomic-embed-text} (via Ollama) to compute 768-dimensional normalized floating-point vectors for session summaries upon completion. These vectors are indexed into a \texttt{vec0} virtual table:

\begin{lstlisting}[language=SQL]
CREATE VIRTUAL TABLE IF NOT EXISTS sessions_vec USING vec0(
    session_id TEXT PRIMARY KEY,
    embedding FLOAT[768] distance_metric=cosine
);
\end{lstlisting}

\section{Reciprocal Rank Fusion (RRF)}
Neither lexical nor semantic search alone suffices for comprehensive memory retrieval. Finch.ai unifies both modalities using \textbf{Reciprocal Rank Fusion (RRF)}.

\subsection{Mathematical Formulation}
Let $D$ denote the corpus of past sessions. Suppose lexical retrieval produces a ranked list $R_{\text{lex}}$ and vector search produces a ranked list $R_{\text{vec}}$. For a session $d \in D$, let $r_{\text{lex}}(d)$ and $r_{\text{vec}}(d)$ represent the 1-based ordinal rank of $d$ within each list (or $\infty$ if unranked).

The composite RRF score is computed as:
\begin{equation}
    \text{RRF\_Score}(d) = \sum_{m \in \{\text{lex}, \text{vec}\}} \frac{1}{k + r_m(d)}
\end{equation}
where $k$ is a smoothing constant that dampens the influence of outlier top-ranked items from any single modality. Finch.ai sets the default:
\begin{equation}
    k = 60
\end{equation}

When candidate documents are scored, results are sorted in descending order of $\text{RRF\_Score}(d)$. Sessions appearing near the top of both indices receive the highest composite weight.

\section{Column-Level Storage Compression (Zstandard)}
To maintain a compact disk footprint over months of heavy usage, Finch.ai incorporates column-level compression (\texttt{ai\_assistant/storage/compression.py}) using Zstandard (\texttt{zstd}) level 3 with automatic \texttt{zlib} fallback.

\begin{figure}[htbp]
\centering
\begin{tikzpicture}[node distance=1.5cm, auto]
    \node (plain) [draw, rectangle, fill=gray!10, rounded corners, text width=3cm, text centered] {Plaintext Payload\\(\texttt{str})};
    \node (check) [draw, diamond, aspect=1.8, fill=yellow!15, text width=2.5cm, text centered, right=of plain] {Length $>$\\1,024 bytes?};
    \node (zstd) [draw, rectangle, fill=blue!10, rounded corners, text width=3.2cm, text centered, right=of check] {Zstandard Compress\\(\texttt{zstd.compress}, lvl 3)};
    \node (tag) [draw, rectangle, fill=green!10, rounded corners, text width=3cm, text centered, below=of zstd] {Magic Header Prepended\\(\texttt{b"\textbackslash x28\textbackslash xb5\textbackslash x2f\textbackslash xfd"})};
    \node (raw) [draw, rectangle, fill=gray!10, rounded corners, text width=3cm, text centered, below=of check] {Store Uncompressed};

    \draw [->, thick] (plain) -- (check);
    \draw [->, thick] (check) -- node[above] {Yes} (zstd);
    \draw [->, thick] (check) -- node[left] {No} (raw);
    \draw [->, thick] (zstd) -- (tag);
\end{tikzpicture}
\caption{Column-Level Message Payload Compression Logic}
\label{fig:column_compression}
\end{figure}

Payloads exceeding 1,024 bytes are compressed prior to database insertion. Decompression is transparent upon retrieval, maintaining zero disruption to querying while slashing database size by up to 90\%.
